%% ============================================================
%% Discussion and Future Work
%% ============================================================

\section{Discussion and Future Work}
\label{sec:discussion}

\subsection{Interpretation of Key Results}

\textbf{Detection latency as the primary metric.}
Single-snapshot metrics (F1, AUC-ROC) are inadequate for NTN trust
evaluation because on-off attackers deliberately alternate good and bad
behaviour within observation windows.  At any snapshot, attackers in
their ``good'' phase are indistinguishable from legitimate nodes —
a fundamental evasion property that all five models experience equally.
Detection latency, which integrates temporal trust evolution, provides
a more informative measure of a framework's real-time response
capability.  On this metric, the proposed framework achieves a
$4.3\times$ and $6.1\times$ improvement over Static~DLT and
Centralized~Trust respectively.

\textbf{False positive rate and threshold calibration.}
With a fixed threshold $\theta=0.40$ and round duration
$T_{\mathrm{round}}=60$\,s, the steady-state trust of legitimate
Ground, LEO, and UAV nodes (Eq.~\ref{eq:steady_state}) falls between
0.264 and 0.356 — below $\theta=0.40$.  This causes
FPR\,=\,26.8\% for the proposed model.
Equation~\ref{eq:theta_opt} provides a closed-form expression for the
segment-optimal threshold; a global $\theta=0.20$ separates legitimate
from malicious nodes for all segment types and reduces FPR
analytically to $\approx 0$.  Deployment systems should use adaptive
threshold calibration based on network telemetry of the legitimate
node trust distribution.

\textbf{Paillier HE overhead feasibility.}
At 755\,ms per aggregation cycle for 50 contributors and 1024-bit
keys, the overhead is within the 60-second round budget for clusters
of up to 50 ground IoT devices reporting to a HAPS aggregator.
For denser deployments, 2048-bit keys or lattice-based HE schemes
(CKKS) reduce or eliminate linear scaling constraints; this is
addressed as future work below.

\subsection{Future Work}

\textbf{1. Adaptive per-segment threshold.}
The closed-form threshold from Eq.~(\ref{eq:theta_opt}) requires
knowledge of $\bar{s}_{\mathrm{legit}}$, which can be estimated
online via a sliding-window median of trust scores from
well-established nodes.  An online EM-based threshold estimator
that updates $\theta_s$ continuously without knowledge of the
attack fraction is a natural extension.

\textbf{2. Lattice-based HE for satellite uplink constraints.}
The GEO uplink budget (typically $<$\,50\,KB per burst) is exceeded
by Paillier ciphertexts at 100+ contributors (25\,KB each).
CKKS or BFV schemes with batching reduce ciphertext sizes by
$32\times$ while maintaining approximate homomorphic aggregation.

\textbf{3. Federated ZTA policy learning.}
The ZTA policy layer currently applies static RBAC rules.
Federated reinforcement learning across NTN segments — where each
satellite or HAPS learns adaptive access rules from local trust
histories without sharing raw telemetry — is a promising direction
for self-evolving security postures.

\textbf{4. Hardware-in-the-loop validation.}
The simulation adopts an Erd\H{o}s-R{\'e}nyi random graph with
empirically-derived delay distributions.  Integration with a
3GPP NR-NTN channel simulator or a hardware testbed (e.g.,
SDR-based LEO emulator) would validate the link-delay model and
provide measured packet-delivery ratios for feedback scoring.

\textbf{5. Formal privacy guarantees under composition.}
The current privacy analysis bounds leakage per aggregation round
using standard DP composition.  The R{\'e}nyi DP accountant and
zero-concentrated DP framework yield tighter cumulative bounds over
$R=500$ rounds and should replace the basic $(\varepsilon,\delta)$
analysis.

\textbf{6. Cross-orbit DLT partitioning.}
The partition experiment (Exp2) simulates a single 50-round partition
event.  GEO–LEO handover creates periodic disconnections every
$\sim$\,90\,minutes; evaluating the reconciliation protocol under
repeated, structured partitions with varying inter-partition intervals
is required before LEO mega-constellation deployment.
